\subsection{2. Bestätigte Mathematik}\label{bestuxe4tigte-mathematik}

Der N=3-Raum zerfällt als Darstellung von Spin(10) × SU(4) in sieben
Isotypien. Drei haben Multiplizität zwei, die anderen vier Multiplizität
eins. Die zwei Fockteile Λ³(C⁶⁴) und C⁶⁰⊗C⁶⁴ sind jeweils
multiplizitätsfrei; der gesamte N=3-Raum ist es ausdrücklich
\textbf{nicht}.

Die irreduziblen Dimensionen sind 24000, 11200, 2688, 2880, 576, 320,
64; die mit Multiplizitäten gewichtete Summe ist 45504. Die
Spin(10)-Darstellung mit Dimension 16 in BF und Λ³ ist in den
angegebenen Gewichten die konjugierte Spinordarstellung; die
Dimensionskurzschreibweise \codeword{16} darf diese Information nicht
ersetzen.

Bestätigt sind:

\begin{itemize}
\tightlist
\item
  die exakte Zerlegung über Gewichtsmultiplizitäten und
  Racah-Alternation;
\item
  \codeword{C(560)=117/4} in der erklärten Casimir-Normierung;
\item
  die Quellenintertwiner für alle 60 Lie-Erzeuger;
\item
  \codeword{rank C3=3776}, \codeword{dim ker C3=37888};
\item
  Spektrum von \codeword{C3 C3†}: 0,7,10,12 mit Multiplizitäten
  64,2880,576,320;
\item
  die drei dunklen Typen mit Gesamtcasimir 45, aber getrennten
  Spin-Casimiren 141/4,117/4,165/4;
\item
  für den \textbf{gewährten} Satz X,Nb eine assoziative Algebra der
  komplexen Dimension 14;
\item
  für alle punktweise symmetrieinvarianten Operatoren die Dimension 16;
\item
  nach zusätzlicher Gewährung aller 60 Gruppenerzeuger die assoziative
  Dimension 743583744 und Kommutantdimension 7.
\end{itemize}

Die Eigenwertzuordnung im Code prüft pro S-Eigenwert einen Vektor, nicht
eine vollständige Basis dieses Eigenraums. Zusammen mit der unabhängig
geprüften Multiplizitätsfreiheit, der Kommutation mit allen Erzeugern
und der eindeutigen Dimensionszuordnung der BF-Teilsummen ist die
Zuordnung trotzdem abgesichert. Die Guard-Beschriftung »entire
eigenspace« ist als direkte Beschreibung dieser einen Rechnung zu weit.

\subsection{3. Invarianz und Kovarianz
auseinanderhalten}\label{invarianz-und-kovarianz-auseinanderhalten}

Mit \codeword{A0=Alg*(X,Nb)} und Gruppenwirkung ρ bezeichne

\codeword{C = End\_G(H3) = ⊕\allowbreak{}\_i End(C\textasciicircum{}\{m\_i\}) ⊗ I\_\{d\_i\}}.

Dann gilt \codeword{A0 ⊂ C}, \codeword{dim A0=14} und
\codeword{dim C=Σ m\_i²=16}. Jedes Wort aus X und Nb kommutiert mit
ρ(G). Nichtzentrale Lie-Erzeuger können deshalb aus diesem Alphabet
nicht entstehen. Dieser Ausschluss ist korrekt und bleibt richtig, wenn
nur ein fixes H statt zweier unabhängiger Kontrollen gewährt ist.

Nach Gewährung der Gruppenerzeuger entsteht dagegen

\codeword{Afull = ⊕\allowbreak{}\_i End(C\textasciicircum{}\{m\_i\} ⊗ V\_i)}.

Diese Algebra ist \textbf{unter Gruppenkonjugation stabil}, aber ihre
Operatoren sind nicht sämtlich invariant. In
\codeword{symmetry\_availability.py} ist die Kennzeichnung
\codeword{symmetry\_invariant: true} für diese Zeile daher falsch. Der
Kommutant von Afull ist die sieben-dimensionale skalare Blockmitte.

Sieben ist ein Boden, solange alle zusätzlichen Operationen diese sieben
Isotypie-Projektoren erhalten; insbesondere senken zusätzliche
punktweise G-invariante Operationen den bereits erreichten Kommutanten
nicht weiter. Sieben ist aber \textbf{kein} universeller Boden für unter
G stabile Algebren oder alle denkbaren nativen Verträge.
\codeword{End(H3)} ist unter jeder Gruppenkonjugation stabil und hat nur
den skalaren Kommutanten. Ein exaktes 2×2-Gegenmodell in der
Zusatzprüfung macht die Unterscheidung direkt sichtbar.

Die Zahlen sind zudem Dimensionen \textbf{komplexer assoziativer
Sternalgebren}, keine nachgewiesenen Dimensionen einer dynamischen
Lie-Algebra oder erreichbaren Unitärgruppe. Burnside liefert nicht
automatisch beliebige ausführbare Gatter auf jedem irreduziblen Block.

\subsection{4. Verfügbarkeit ist keine universelle
Ja/Nein-Frage}\label{verfuxfcgbarkeit-ist-keine-universelle-janein-frage}

Ein vorgegebener Hamiltonoperator allein gewährt nicht bereits das
unabhängige Schalten von X und Nb. Am erklärten Punkt g/Δ=1/20 besitzt
ein fixes H im N=3-Raum acht verschiedene Energien, also nur eine
acht-dimensionale kommutative Spektralalgebra. Die großzügigere Annahme
unabhängiger Kontrollen X,Nb ergibt 14.

Eine treue, auf die Multiplizitätsräume reduzierte rationale
Matrixrechnung ergibt:

{\def\LTcaptype{none} % do not increment counter
\begin{longtable}[]{@{}lr@{}}
\toprule\noalign{}
Zusätzlich gewährter Satz & Assoziative Dimension \\
\midrule\noalign{}
\endhead
\bottomrule\noalign{}
\endlastfoot
Nur fixes H=Nb+X/20 & 8 \\
X,Nb & 14 \\
X,Nb plus Projektor auf genau einen dunklen Typ & 15 \\
X,Nb plus Spin(10)-Casimir & 16 \\
X,Nb plus SU(4)-Casimir & 16 \\
X,Nb plus Gesamtcasimir & 14 \\
\end{longtable}
}

Bereits \textbf{einer} der getrennten Casimire reicht zur vollständigen
Trennung der drei dunklen Typen. Die zwei fehlenden Algebradimensionen
sind kein Beweis, dass genau zwei physische Griffe fehlen. Die Dimension
15 zeigt eine mögliche invariante Zwischenstufe. Die behauptete
universelle Dichotomie »getrennte Griffe ja oder unverändert 14« und das
\codeword{IF AND ONLY IF} zur Verfügbarkeit der vollen Symmetrie sind
daher zu stark.

Die zusätzlichen Projektoren und Casimirkontrollen werden hier nur als
algebraische Gegenbeispiele verwendet. Ihre Implementierung aus dem
Compiler wird nicht behauptet.

\subsection{5. Der Clock-Satz ist im gelieferten Checker nicht
geprüft}\label{der-clock-satz-ist-im-gelieferten-checker-nicht-gepruxfcft}

\codeword{operation\_symmetry.py:478} schreibt den Normalisator- und
Nichts-hinzufügen-Satz als \codeword{need(True, ...)}. Insgesamt enthält
diese Datei neun solche als Prüfbedingungen gezählten theoretischen
Folgerungen. Einige sind durch die vorausgehende Mathematik gut
begründet; der bloße grüne Guard zertifiziert sie jedoch nicht
unabhängig.

Die direkte Clock-Definition liegt in
\codeword{universalraum-native-operations-ground-response-20260915/common.py},
Funktion \codeword{clock\_lift}, ab Zeile 187: eine Permutation der fünf
Oszillatorachsen wird mit Exterior-Vorzeichen auf den geraden Spinorraum
gehoben; \codeword{GF=G16⊗I4}. Der Bosonlift wird entsprechend gebaut
und die Tensorintertwining-Gleichung geprüft. Eine neue direkte Prüfung
der Clock-/Casimir-Beziehungen erfolgt im parallelen Hauptstrang, nicht
in diesem Audit.

\subsection{6. Feldtyp: was tatsächlich erzwungen
ist}\label{feldtyp-was-tatsuxe4chlich-erzwungen-ist}

Für ein \textbf{lokales, ableitungsfreies Bilinear zweier gleichhändiger
Weylfelder}, das alle 64 inneren Quellenmarken unverändert als
unabhängige innere Komponenten und genau den gegebenen antisymmetrischen
Tensor W verwendet, gilt

\codeword{(1/2,0) ⊗ (1/2,0) = (1,0) ⊕\allowbreak{} (0,0)}.

Die skalare Epsilon-Kontraktion zusammen mit antisymmetrischem W ist im
gemeinsamen Index symmetrisch und verschwindet wegen
Grassmann-Antikommutation. Der symmetrische Spinortensor, also der
(1,0)-Kanal, bleibt nichtverschwindend. Im genannten Vertrag ist dies
korrekt. Ein zusätzliches unabhängiges gleichgeladenes Hilfsdublett
erlaubt stattdessen wieder die skalare Kontraktion.

Die Zahlen 180/128 beziehungsweise 60/256 zählen \textbf{komplexe lokale
Feldkomponenten im jeweiligen direkten Tensorprodukt-Ansatz}. Sie zählen
nicht automatisch zusätzliche unabhängige physische Oszillatoren oder
propagierende Freiheitsgrade. Diese hängen von Kinetik,
Nebenbedingungen, positiver Energie, Teilchen-/Antiteilchenstruktur und
dem tatsächlichen Feldadapter ab. Deshalb folgt aus der Multiplikation
mit zwei oder drei kein allgemeiner Ausschluss jeder relativistischen
Lesart mit den nativen Modenzahlen.

Vor allem folgt daraus keine vorgeschriebene Reihenfolge »erst räumliche
Skalierung, dann Feldwörterbuch«. Lokale Darstellungstypen und mögliche
Kinetik können vor einer Skalierungsrechnung geprüft werden. Welche
Feldkomponenten tatsächlich aus dem räumlichen Quellprozess entstehen,
muss anschließend gemeinsam mit dem Adapter untersucht werden.

\subsubsection{Konkreter Gegencheck jenseits der ableitungsfreien
Klasse}\label{konkreter-gegencheck-jenseits-der-ableitungsfreien-klasse}

Der allgemeine Satz »ein Vektor verlangt ψ†ψ und scheitert deshalb an
der Ladung« ist falsch, sobald Ableitungen zugelassen werden.
Beispielsweise ist

\codeword{J\textasciicircum{}A\_mu = Σ\_IJab M\textasciicircum{}A\_IJ ε\_ab ψ\_Ia ↔∂\_mu ψ\_Jb}

ein Lorentzvektor mit Ladung −2. Die Ableitungsantisymmetrisierung macht
ihn bei antisymmetrischem M nichtverschwindend. Für M=ε und zwei inneren
Marken liefert die exakte Grassmann-Jetrechnung vier nichtverschwindende
Monome mit Koeffizienten +2,−2,−2,+2. Ein entsprechend geladener
Vektorvermittler könnte algebraisch durch
\codeword{b†\_mu J\textasciicircum{}mu + h.c.} koppeln, ohne
Ladungsverletzung.

Das ist \textbf{nur ein Gegenbeispiel gegen den überbreiten Ausschluss},
keine native Lösung: Es fügt eine Ableitung und einen Feldtyp hinzu; bei
kanonischen 4D-Felddimensionen ist es ein höherdimensionaler
Kopplungsterm. Gesunde Kinetik, Quellenherkunft, Skalierung und T1-T8
werden dadurch nicht bewiesen.

\subsection{7. Konsequenz}\label{konsequenz}

Die neuen Quellen verkleinern und präzisieren den endlichen
N=3-Operationsraum. Sie schließen weder die physische Verfügbarkeit der
Operationen noch den Lorentzadapter. Die sachlich tragfähige Integration
übernimmt die Zerlegung, Casimire, 14/16/7 im jeweils erklärten Vertrag
sowie den eingeschränkten Feldtypsatz; sie ersetzt die überbreiten
Ausschlüsse und die behauptete Programmumkehr durch explizite
Bedingungen.

\begin{center}\rule{0.5\linewidth}{0.5pt}\end{center}

\section{Anhang B: unabhängige Krylov- und
Feldprüfung}\label{anhang-b-unabhuxe4ngige-krylov--und-feldpruxfcfung}

\section{Unabhängige Prüfung: Krylov-Verzweigung und
Feldwörterbuch}\label{unabhuxe4ngige-pruxfcfung-krylov-verzweigung-und-feldwuxf6rterbuch}

Arbeitsstand 2026-09-15. Geprüft wird die Anlage
\codeword{7380c383-f364-4ef8-8066-24aa4d180d85} gegen den Vertrag
\codeword{universalraum-native-operations-ground-response-20260915}.
Fremde Quellen und Ausgaben wurden nicht verändert. Die hier angegebenen
Matrix- und Graphbefunde beziehen sich auf denselben gepinnten Tensor W
mit SHA-256
\codeword{3f00a0892157a691e4ef26920c3702772c7bf76a250fad98ff04a079bb64b763}.

\subsection{1. Der neue Nichtschluss ist exakt
bestätigt}\label{der-neue-nichtschluss-ist-exakt-bestuxe4tigt}

Für \codeword{v\_n=(T†)\textasciicircum{}n F} habe ich
\codeword{u=T T† v\_2=T v\_3} unmittelbar aus den 480 nativen Paaren neu
kontrahiert. Das materialisiert nicht die 15.252.960 Komponenten von v₃.
Es arbeitet mit 108.240 Komponenten von v₂ und 293.280
nichtverschwindenden Komponenten des Ergebnisses u. Es wurden 45.771.360
zulässige Erzeugungsübergänge und 139.289.760 anschließende
Entnahmebeiträge exakt mit ganzzahliger Fermion-Parität und
Boson-Multiplizität zusammengeführt.

Die unabhängige Rechnung liefert

\[\adjustbox{max width=\linewidth}{$\displaystyle 
\|v_2\|^2=439680,\quad
\langle v_2,u\rangle=575078400,\quad
\|u\|^2=752194252800.
$}\]

Damit folgen rein rational

\[\adjustbox{max width=\linewidth}{$\displaystyle 
\alpha=\frac{299520}{229},\qquad
w_2=u-\alpha v_2,\qquad
\|w_2\|^2=\frac{5001523200}{229}>0,
$}\]

und

\[\adjustbox{max width=\linewidth}{$\displaystyle 
\frac{\|w_2\|^2}{\nu_3}=\frac{1809}{47632}.
$}\]

Der Ein-Vektor-pro-Bosonlage-Ansatz ist somit nicht invariant. Da T die
native Symmetrie erhält und v₂ und u Singuletts sind, liefert w₂ einen
zweiten unabhängigen Singulettvektor in derselben Bosonlage. Das ist ein
Strukturresultat, kein physikalischer Mehrteilchen- oder
Raumzeitnachweis.

\subsection{2. Das mitgelieferte Ground-JSON ist veraltet und teilweise
falsch}\label{das-mitgelieferte-ground-json-ist-veraltet-und-teilweise-falsch}

Die aktuelle Pythonquelle hat SHA-256
\codeword{8ae0f2f103ba0cb0a5cb832681a33130b6756fe627395368b6be635587a2eb02}.
Das vorhandene \codeword{native\_ground\_state.json} nennt dagegen
\codeword{046ffa3f41d9dd365e2c9a199e938734bf735d762f23373a3f85b541c3da30fe}
und enthält noch

\[\adjustbox{max width=\linewidth}{$\displaystyle 
\|w_2\|^2_{\rm alt}=1145348812800
=229^2\|w_2\|^2.
$}\]

Sein \codeword{closure\_defect=1.5226768996658837} ist schon als
relative orthogonale Norm unmöglich. Das daraus abgeleitete
Ritz-Ergebnis −1,18826198 ist ungültig. Der spätere Anlagentext
korrigiert den Normfehler, aber die JSON-Belegkette wurde noch nicht
entsprechend neu erzeugt. Das grüne \codeword{status: PASS} dieser
historischen Ausgabe darf deshalb nicht mit dem korrigierten Text
zusammengezogen werden.

Die korrigierte Norm und die folgenden kleinen Matrizen wurden hier
unabhängig reproduziert. Die ganze fremde große Ausführung wurde
ausdrücklich nicht erneut gestartet.

\subsection{3. Die sechste Dezimalstelle ist richtig - aber nur für die
eine zusätzliche
Richtung}\label{die-sechste-dezimalstelle-ist-richtig---aber-nur-fuxfcr-die-eine-zusuxe4tzliche-richtung}

Die neue numerische Diagonalisierung ergibt:

{\def\LTcaptype{none} % do not increment counter
\begin{longtable}[]{@{}
  >{\raggedright\arraybackslash}p{(\linewidth - 6\tabcolsep) * \real{0.2000}}
  >{\raggedleft\arraybackslash}p{(\linewidth - 6\tabcolsep) * \real{0.2667}}
  >{\raggedleft\arraybackslash}p{(\linewidth - 6\tabcolsep) * \real{0.2667}}
  >{\raggedleft\arraybackslash}p{(\linewidth - 6\tabcolsep) * \real{0.2667}}@{}}
\toprule\noalign{}
\begin{minipage}[b]{\linewidth}\raggedright
g/Δ
\end{minipage} & \begin{minipage}[b]{\linewidth}\raggedleft
vierdimensionale Ritz-Matrix
\end{minipage} & \begin{minipage}[b]{\linewidth}\raggedleft
plus w₂
\end{minipage} & \begin{minipage}[b]{\linewidth}\raggedleft
Änderung
\end{minipage} \\
\midrule\noalign{}
\endhead
\bottomrule\noalign{}
\endlastfoot
1/20 & −1,0942308394409612 & −1,0942318710142458 & −1,0315733·10⁻⁶ \\
1/40 & −0,29535718178032033 & −0,29535720496683315 & −2,3186513·10⁻⁸ \\
1/100 & −0,04789424798601009 & −0,04789424801311775 &
−2,7107663·10⁻¹¹ \\
\end{longtable}
}

Das sind numerische Eigenwerte exakter kleiner Ritz-Matrizen, keine nach
außen gerundeten vollständigen Spektraleinschließungen. Durch
Eliminieren der einzigen neuen w₂-Koordinate bekommt die v₃-Diagonale
die exakt bestimmte energieabhängige Korrektur

\[\adjustbox{max width=\linewidth}{$\displaystyle 
-\frac{g^2(1809/47632)}{2\Delta-E}.
$}\]

Das erklärt die kleine Änderung durch diesen einen Zweig. Es sagt nichts
über die ausgelassenen höheren Bosonlagen.

\subsubsection{Eigene einfache Folgeherleitung: Das vollständige
Ritz-Residuum}\label{eigene-einfache-folgeherleitung-das-vollstuxe4ndige-ritz-residuum}

Sei ψ₃ der normierte Ritz-Eigenvektor in \codeword{span(v₀,…,v₃)}, und
c₃ seine Komponente entlang \codeword{v₃/√ν₃}. Außerhalb dieses Raums
führt Hψ₃ genau in zwei orthogonale Richtungen: v₄ und w₂. Deshalb gilt
exakt

\[\adjustbox{max width=\linewidth}{$\displaystyle 
\|(H-E_{\rm Ritz})\psi_3\|^2
=g^2|c_3|^2\frac{\nu_4+\|w_2\|^2}{\nu_3}.
$}\]

Hier wird \codeword{ν₄=952296652800} aus dem früher vollständig
geprüften Momentbeleg übernommen; die vierte Ordnung wurde in dieser
Runde nicht erneut enumeriert.

Bei g/Δ=1/20 folgt numerisch ein Residuum von rund \textbf{0,373063 Δ}.
Der w₂-Zweig trägt weniger als \textbf{23 Millionstel} des quadrierten
Residuums. Die große ausgelassene Richtung ist v₄, nicht w₂. Aus dem
winzigen Nichtschlussanteil darf daher keine praktisch vollständige
Konvergenz abgeleitet werden.

Außerdem liegt die neue Fünf-Zustands-Ritzenergie mindestens
\textbf{0,0354 Δ über} der bereits bewiesenen Obergrenze des wahren
Grundzustands. Das ist eine nachweislich noch relevante Lücke,
unabhängig von einer Konvergenzschätzung.

